\begin{description}
\item[(a)] To measure instrumental magnitude we should choose an aperture.
Careful investigation of the image shows that a $5 \times 5$ pixel aperture
is enough to measure $m_I$ for all stars. $m_I$ can be calculated using:
\[ m_I = -2.5\log\!\left(\frac{\sum_{i=1}^{N} I_{i(\mathrm{star})} - N\bar{I}_{\mathrm{Sky}}}{\mathrm{Exp}}\right) \]
where $I_{i(\mathrm{star})}$ is the pixel value for each pixel inside the
aperture, $N$ is number of pixels inside the aperture,
$\bar{I}_{\mathrm{Sky}}$ is the average of sky value per pixel taken from
dark part of image and Exp is the exposure time. Table (1.4) lists values
for $m_I$ and $Zmag$ calculated for all three identified stars.
\[ \bar{I}_{\mathrm{Sky}} = 4.42 \]
\[ N = 25 \]
\[ Exp = 450 \]

\begin{center}
Table (1.4)\\[2pt]
\begin{tabular}{|c|c|c|c|}
\hline
Star & $m_I$ & $m_t$ & $Zmag$ \\
\hline
1 & $-3.02$ & 9.03 & 12.38 \\
\hline
3 & $-5.85$ & 6.22 & 12.40 \\
\hline
4 & $-4.04$ & 8.02 & 12.39 \\
\hline
\end{tabular}
\end{center}

\item[(b)] Average $Zmag = 12.4$.

\item[(c)] Following part (a) for stars 2 and 5, we can calculate true
magnitudes ($m_t$) for these stars.

\begin{center}
Table (1.5)\\[2pt]
\begin{tabular}{|c|c|c|}
\hline
Star & $m_I$ & $m_t$ \\
\hline
2 & $-2.13$ & 9.93 \\
\hline
5 & $-0.66$ & 11.4 \\
\hline
\end{tabular}
\end{center}

\item[(d)] Pixel scale for this CCD is calculated as
\[ p = \frac{\text{pixel size}}{\text{focal length}} \times \frac{180 \times 3600}{\pi} = 4.30'' \]

\item[(e)] Average sky brightness:
\[ m_{sky} = -2.5\log\frac{\bar{I}_{\mathrm{Sky}}}{(Exp)(p)^2} + Zmag = 20.6 \]

\item[(f)] To estimate astronomical seeing, first we plot pixel values in
$x$ or $y$ direction for one of the bright stars in the image. As plot (1)
shows, the FWHM of pixel values which is plotted for star 3, is 1 pixel,
hence astronomical seeing is equal to
\[ seeing \cong 4'' \]

% FIGURE NEEDED: Plot (1) -- pixel-value profile of star 3 (Gaussian-like
% curve peaking near 45000 at pixel 4, with FWHM marked around 1 pixel)
% (2009_da_sol.pdf page 3)
\end{description}
